\documentclass[11pt,a4paper]{article}
\usepackage[T1]{fontenc}
\usepackage[utf8]{inputenc}
\usepackage[margin=25mm]{geometry}
\usepackage{amsmath,amssymb,amsthm,booktabs}
\usepackage[hidelinks]{hyperref}
\newtheorem{proposition}{Proposition}
\newtheorem{lemma}{Lemma}
\newcommand{\norm}[1]{\left\lVert#1\right\rVert}
\title{Adam under Orthogonal Coordinate Changes\\
\large Module 2: A Positive-Epsilon Counterexample and Fixed-Gradient Replay}
\author{AISC5020 Project Team\\Module 2}
\date{18 October 2026}
\begin{document}
\maketitle

\begin{abstract}
We examine how Adam's coordinatewise moment normalization responds to
orthogonal changes of coordinates. A one-step counterexample establishes
failure of general rotation equivariance for every finite positive scalar
epsilon, without replacing the implemented update by a zero-epsilon limit.
Signed permutations remain an exact special case. We then distinguish a
fixed-gradient replay diagnostic from a full optimization trajectory: the
former holds the input history fixed up to rotation, whereas the latter
includes feedback from diverging parameter positions. A reproduced
two-dimensional diagnostic contains 350 rows over seven angles and 50 updates.
It tests the mechanism, not optimizer rankings or convergence rates.
\end{abstract}

\section{Problem and update convention}
Let $D$ be positive diagonal and $Q$ orthogonal. The paired problems are
\[
 f(x)=\tfrac12x^\top Dx,\qquad
 \widetilde f(y)=\tfrac12y^\top(QDQ^\top)y,\qquad y_0=Qx_0.
\]
They have the same eigenvalues, spectral condition number, minimum value,
and corresponding objective geometry. This does not imply that every
optimization algorithm produces corresponding trajectories.

We use bias-corrected Adam~\cite{adam}, with zero initial moments and a scalar
$\epsilon>0$ outside the square root. At update $t$, the gradient is evaluated
before moving, $g_t=\nabla f(x_{t-1})$, and
\begin{align}
 m_t&=\beta_1m_{t-1}+(1-\beta_1)g_t,
 &v_t&=\beta_2v_{t-1}+(1-\beta_2)(g_t\odot g_t),\label{eq:moments}\\
 \widehat m_t&=\frac{m_t}{1-\beta_1^t},
 &\widehat v_t&=\frac{v_t}{1-\beta_2^t},\\
 A_t&=\frac{\widehat m_t}{\sqrt{\widehat v_t}+\epsilon\mathbf1},
 &x_t&=x_{t-1}-\eta A_t.\label{eq:direction}
\end{align}
Here $\beta_1,\beta_2\in[0,1)$ and $\eta>0$. Squares, square roots and
quotients of vectors are componentwise. The direction $A_t$ has not yet been
multiplied by the learning rate. In code, \texttt{adam\_direction} updates the
moments and time counter and returns one direction vector;
\texttt{adam\_step} additionally moves the parameter position.

For a supplied gradient history, write $A_t=A(g_1,\ldots,g_t)$.
Rotation equivariance would require
\begin{equation}
 A(Qg_1,\ldots,Qg_t)=QA(g_1,\ldots,g_t)
 \label{eq:equivariance}
\end{equation}
for every history and orthogonal $Q$, using the same scalar parameters.
The following example disproves that general statement.

\section{A hand-checkable counterexample with positive epsilon}
\begin{proposition}[P3: general rotation equivariance fails]
For every finite $\epsilon>0$ and $\beta_1,\beta_2\in[0,1)$, there is a
positive-definite quadratic and a paired orthogonal rotation for which
Adam's first directions fail \eqref{eq:equivariance}.
\end{proposition}
\begin{proof}
Take $D=\operatorname{diag}(1,2)$, $x_0=(1,0)^\top$, and let $Q$ rotate
through $\pi/4$. Write $a=1/\sqrt2$. The initial gradient and its rotated
counterpart are $g_1=(1,0)^\top$ and $Qg_1=(a,a)^\top$.
Starting from zero moments, bias correction gives
$\widehat m_1=g_1$ and $\widehat v_1=g_1\odot g_1$, independently of the
two moment coefficients. Therefore
\[
 A(Qg_1)=\frac{(a,a)^\top}{a+\epsilon},\qquad
 QA(g_1)=\frac{(a,a)^\top}{1+\epsilon}.
\]
The denominators differ for every finite positive $\epsilon$. More explicitly,
\begin{equation}
 \norm{A(Qg_1)-QA(g_1)}_2
 =\frac1{a+\epsilon}-\frac1{1+\epsilon}
 =\frac{1-a}{(a+\epsilon)(1+\epsilon)}>0.
 \label{eq:counterexample}
\end{equation}
Both problems have corresponding initial gradients, so the first parameter
discrepancy is $\eta$ times this nonzero direction discrepancy.
\end{proof}

For $\epsilon=10^{-8}$, the absolute direction discrepancy in
\eqref{eq:counterexample} is approximately $0.414214$.
This is a separate analytical example, not an extra sample in the
350-row experiment below. It establishes existence of a failure; it does
not claim that every rotation or every initial point causes a discrepancy.

\begin{lemma}[Signed permutations remain equivariant]
If $S$ is a signed permutation matrix, $P=|S|$ its unsigned permutation,
and the two supplied histories satisfy $g'_t=Sg_t$, then Adam with zero
initial moments and shared scalar parameters satisfies
$m'_t=Sm_t$, $v'_t=Pv_t$, and $A'_t=SA_t$ in exact arithmetic.
\end{lemma}
\begin{proof}
The claims hold at initialization. The first-moment update is linear, and
$(Sg)\odot(Sg)=P(g\odot g)$, so induction proves the two moment identities.
Scalar bias correction preserves them. Since $P\mathbf1=\mathbf1$ and
$\sqrt{Pv}=P\sqrt v$, the denominator is permuted without signs while the
numerator has the appropriate signs. Componentwise division therefore gives
$A'_t=SA_t$.
\end{proof}
For paired quadratic problems, this direction identity and induction on
positions also imply $y_t=Sx_t$. The identity and a $90^\circ$ rotation in
two dimensions are such controls. Notice that second moments use $P$, not
$S$: directly asserting $v'_t=Sv_t$ would be incorrect. The lemma also does
not extend to arbitrary rotations within a repeated eigenspace. Module 5
addresses Adam in a specified fixed eigenbasis under separate assumptions.

\section{Two diagnostics with different meanings}
\subsection{Fixed-gradient replay}
We first generate one canonical GD gradient history, with its own learning
rate $\alpha$:
\[
 g_t=Dx_{t-1}^{\rm GD},\qquad
 x_t^{\rm GD}=x_{t-1}^{\rm GD}-\alpha g_t,
 \qquad
 g_t=D(I-\alpha D)^{t-1}x_0.
\]
GD commutes with orthogonal rotation, so rotating this history gives the
corresponding history $Qg_t$ without another optimization run. If gradients
are stored as rows of a matrix $G$, the rotated row history is $GQ^\top$,
because $(Qg_t)^\top=g_t^\top Q^\top$.

\texttt{replay\_directions} creates a fresh zero moment state and processes
the supplied gradients sequentially. It returns a $T\times p$ direction
matrix, without moving parameters or generating new gradients. Rotation is
performed by the matrix multiplication before this function is called.
Two independent replays yield $A_t$ and $A'_t$, and we report
\begin{equation}
 n_t=\norm{A'_t-QA_t}_2,\qquad
 b_t=\norm{A_t}_2,\qquad
 \delta_t=\frac{n_t}{b_t+10^{-12}}.
 \label{eq:defect}
\end{equation}
The $10^{-12}$ term protects this diagnostic denominator; it is distinct
from Adam's optimizer epsilon. The normalized direction defect is neither
a loss difference nor a convergence rate. Positive epsilon and this
protection term also prevent an unconditional claim of scale invariance.

\subsection{Actual Adam trajectories}
We separately run ordinary Adam on both paired quadratics. Each run computes
new gradients at its actual positions. The stored paths contain the initial
point and all updated positions, giving $T+1$ rows; a gradient history instead
has $T$ rows and begins with the gradient at the initial point.
The position discrepancy is
\begin{equation}
 e_t=\frac{\norm{y_t-Qx_t}_2}{\max(1,\norm{x_0}_2)}.
 \label{eq:position}
\end{equation}
Once the paths separate, later gradients need not correspond. This diagnostic
therefore includes feedback from parameter movement, whereas replay isolates
the update response to histories that are exactly paired by construction.
Neither metric alone ranks the quality of the final optimized solutions.

\section{Reproduced module diagnostic}
The default configuration uses seed 502002,
$D=\operatorname{diag}(0.1,1)$, initial objective 1,
$\alpha=0.5$, $\eta=0.1$, $(\beta_1,\beta_2)=(0.9,0.999)$,
$\epsilon=10^{-8}$, and 50 updates. Seven angles from $0^\circ$ to
$90^\circ$ in $15^\circ$ increments give 350 diagnostic rows. Checkpoints
are updates 1, 2, 10 and 50. The canonical initial point is approximately
$(3.85243548,0.71824375)^\top$ and is shared across all angles through
rotation. Reproduction used Python 3.12.14 and NumPy 2.5.4.

\begin{table}[htbp]
\centering
\caption{Observed direction and position discrepancies. Maxima are over
50 updates; the final position column is at update 50.}
\label{tab:results}
\begin{tabular}{rrrr}
\toprule
Angle & Maximum $\delta_t$ & Maximum $e_t$ & Final $e_t$\\
\midrule
$0^\circ$  & 0 & 0 & 0\\
$15^\circ$ & 0.261052 & 0.128376 & 0.0731763\\
$30^\circ$ & 0.9999997 & 0.173925 & 0.104984\\
$45^\circ$ & 0.765367 & 0.112466 & 0.0237806\\
$60^\circ$ & 0.517638 & 0.111343 & 0.0400905\\
$75^\circ$ & 0.261052 & 0.0737234 & 0.0424682\\
$90^\circ$ & $2.174\,10^{-16}$ & $1.417\,10^{-17}$ & 0\\
\bottomrule
\end{tabular}
\end{table}

The identity control is exactly zero; the $90^\circ$ control agrees to
floating-point precision. Intermediate rotations give nonzero replay and
position discrepancies in this fixed experiment. Their magnitudes need not
increase monotonically with angle: they depend on the history's coordinate
components, not just the geometric angle. In particular, the largest replay
defect in this run occurs at $30^\circ$, not $45^\circ$.

A separate local audit reconstructed Adam moments from explicit weighted
gradient sums and independently reconstructed the actual position updates.
It checked five numerical metrics per row, 1,750 scalar comparisons in total;
the largest absolute difference was $1.2213\times10^{-15}$. This corroborates
the default diagnostic's numbers. It is not a convergence proof, an
independent statistical replication, or evidence that all possible
configurations have been validated.

The scope is intentionally limited: one spectrum, one initialization, one
step-size choice and a short fixed horizon. The seven angles share an input
history, so they are not seven independent random trials. No tuning, success
rate estimate, uncertainty interval or wall-clock superiority is inferred.
The recorded protocol is not the frozen team main-study protocol.

\section{Reproduction and integration}
From the repository root in the installed project environment, run
\begin{verbatim}
python -m unittest discover -s tests -p test_adam_diagnostics.py -v
python -m coordinate_geometry.adam_diagnostics \
  --config configs/adam_diagnostics.json \
  --output outputs/adam_diagnostics
\end{verbatim}
Adding \texttt{--plot} generates the figure and requires the optional plot
dependencies. Artifacts include \path{diagnostics.csv}, \path{problem.json},
\path{run_config.json}, \path{summary.json}, and \path{run.log} in the selected
output directory. Rerunning overwrites that directory's corresponding files;
choose a new directory when preserving an earlier run.

The group can extract the theory, method and appropriately scoped results
from this standalone note for its final report. Module 2 contributes a
positive-epsilon counterexample, exact signed-permutation controls, and a
diagnostic separating fixed-input normalization from trajectory feedback.
The group must still agree on parameter budgets, spectral and initialization
sampling, and statistical analysis before drawing its main experimental
conclusions. In particular, recovering correspondence in a common eigenbasis
does not make ordinary ambient-coordinate Adam generally rotation equivariant.

\begin{thebibliography}{1}
\bibitem{adam}
D. P. Kingma and J. Ba.
\newblock Adam: A Method for Stochastic Optimization.
\newblock \emph{International Conference on Learning Representations}, 2015.
\newblock \url{https://arxiv.org/abs/1412.6980}.
\end{thebibliography}
\end{document}
